% Apêndices embutidos — sem referência a documentação externa

\section{Contrato de árvore sintática abstrata e mensagens REST}
\label{app:ast-contract}

% Conteúdo embutido: contratos JSON e comunicação entre microsserviços

\subsection{Contrato de representação da árvore sintática abstrata}

A comunicação entre os microsserviços do pipeline restringe-se a mensagens no formato \emph{JavaScript Object Notation} (JSON). O documento raiz deve declarar obrigatoriamente o campo \texttt{type} com valor \texttt{Program}, seguido do vetor \texttt{declarations}, que agrega declarações de classes, funções, variáveis globais e canais. A serialização produzida pelo analisador sintático utiliza o módulo \texttt{ast\_json.to\_dict}; o analisador semântico e os geradores de código consomem a estrutura por meio de \texttt{ast\_from\_json.from\_dict}.

\begin{table}[htbp]
\centering
\small
\caption{Nós principais do contrato de árvore sintática abstrata em JSON}
\label{tab:ast-contract-nodes}
\begin{tabularx}{\textwidth}{|>{\raggedright\arraybackslash}p{3.2cm}|X|}
\hline
\textbf{Tipo} & \textbf{Campos e semântica} \\
\hline
\texttt{Program} & Raiz do documento; contém \texttt{declarations[]} \\
\texttt{ClassDecl} & \texttt{name}, \texttt{extends} (opcional), \texttt{members[]} \\
\texttt{MethodDecl} & \texttt{name}, \texttt{returnType}, \texttt{parameters}, \texttt{body} \\
\texttt{FuncDecl} & Função global com tipo de retorno e corpo \\
\texttt{ParBlock} / \texttt{SeqBlock} & Blocos de paralelismo e sequência; \texttt{statements[]} \\
\texttt{SendStmt} / \texttt{ReceiveStmt} & Comunicação por canal; identificador do canal e expressão \\
\texttt{ChannelDecl} & Declaração \texttt{s\_channel} ou \texttt{c\_channel} \\
\texttt{NewInstance} & Instanciação \texttt{new}; \texttt{className}, \texttt{arguments[]} \\
\texttt{MethodCall} & Chamada de método; receptor, nome e argumentos \\
\hline
\end{tabularx}
\end{table}

\begin{lstlisting}[caption={Exemplo mínimo de documento JSON para uma classe com herança}]
{
  "type": "Program",
  "declarations": [{
    "type": "ClassDecl",
    "name": "Dog",
    "extends": "Animal",
    "members": [{
      "type": "MethodDecl",
      "name": "bark",
      "returnType": "void",
      "parameters": [],
      "body": { "type": "Block", "statements": [] }
    }]
  }]
}
\end{lstlisting}

\subsection{Contrato de resposta dos geradores de código}

O resultado de qualquer variante de tradução materializa-se na estrutura \texttt{TranslationResult}, composta pelos campos \texttt{output} (texto emitido na execução), \texttt{code} (artefato gerado, quando aplicável) e \texttt{exit\_code} (código de retorno do processo externo de compilação ou interpretação). O acoplamento entre microsserviços mantém-se reduzido: o analisador semântico não conhece o back-end de destino, e os geradores de código não reimplementam a análise léxica ou sintática.

\begin{lstlisting}[caption={Corpo típico de requisição ao gateway central}]
POST /api/v1/process
Content-Type: application/json

{
  "sourceCode": "class Main { void run() { println(\"ok\"); } }",
  "targetVariability": "INTERPRETER",
  "executionMode": "LOCAL"
}
\end{lstlisting}


\section{Processo de instanciação de nova variante de produto}
\label{app:creating-app}

A criação de uma nova configuração de produto sobre o \emph{MiniPar Framework} pressupõe a reutilização integral do núcleo compartilhado (\texttt{minipar-core}, gateway central, microsserviços de análise léxica, sintática e semântica, bem como a orquestração por contêineres). O desenvolvedor da instância implementa exclusivamente os pontos de adaptação previstos pelo padrão \emph{Template Method}, materializados nos métodos \texttt{emit} e \texttt{finalize} da classe derivada de \texttt{AbstractBackendTranslator}. Em seguida, encapsula-se a lógica em um microsserviço autônomo que expõe o endpoint REST de geração ou execução, no mesmo molde do serviço de geração de código em linguagem C. O registro da nova variante no vetor \texttt{BACKEND\_REGISTRY}, a definição da variável de ambiente correspondente e a inclusão do serviço no arquivo de composição de contêineres completam a configuração em tempo de implantação. Por fim, a interface de seleção de variabilidade deve contemplar o novo valor de \texttt{targetVariability}, de modo que o usuário final possa escolher a variante em tempo de execução sem alteração do algoritmo de orquestração.

As instâncias de referência atualmente disponíveis compreendem o interpretador direto sobre a árvore sintática abstrata, o compilador com saída em linguagem C integrado ao compilador \texttt{gcc} com otimização de segundo nível, e a extensão experimental com geração de código Python, esta última destinada a demonstrar a extensibilidade do framework sem modificação dos trechos invariantes.

\section{Protocolo de validação e evidências de execução}
\label{app:e2e}

A validação de ponta a ponta do sistema executou-se mediante o script automatizado \texttt{validate-all.sh}, com a stack completa em execução via \emph{Docker Compose}. Em quinze de junho de 2026, obteve-se aprovação em dezesseis verificações, entre as quais se incluem a saída do interpretador para o exemplo \texttt{08\_interpreter\_ok}, a detecção de erro sintático no exemplo \texttt{05\_parse\_extends\_missing}, a compilação com \texttt{rustc} para o exemplo \texttt{12\_codegen\_rust\_stub}, a geração do tapete de Sierpinski no exemplo \texttt{13\_sierpinski}, a coordenação distribuída nos modos legado e por menu MiniPar (\texttt{14\_distributed\_menu}), a verificação semântica de herança inválida, a geração de código C e Python, o teste de canais por soquete TCP no exemplo \texttt{15\_channels}, e a consulta aos endpoints de variantes, recomendações e saúde agregada dos microsserviços. As saídas textuais dos testes distribuído e de canais encontram-se nas Listagens~\ref{lst:evidence-distributed} e~\ref{lst:evidence-channels}.

\section{Glossário}
\label{app:glossary}

\begin{description}[style=nextline]
    \item[Trecho invariante (\emph{frozen-spot})] Segmento do framework que o desenvolvedor da instância não substitui, como o método \texttt{translate} e a orquestração do gateway.
    \item[Ponto de adaptação (\emph{hotspot})] Segmento variável preenchido na instância, em especial \texttt{emit} e \texttt{finalize}.
    \item[Instância de aplicação] Configuração da linha de produto sobre o chassi comum, compreendendo contêineres, pontos de adaptação implementados e registro de variantes.
    \item[Variabilidade] Dimensão configurável da linha de produto, como modo de execução, back-end de tradução e ambiente local ou distribuído.
    \item[\emph{Asset}] Microsserviço ou módulo reutilizável em múltiplas configurações de produto.
    \item[\emph{Binding time}] Momento em que uma variante é vinculada, seja em tempo de implementação, implantação ou execução.
    \item[Princípio de Hollywood] Forma de inversão de controle na qual o framework invoca os pontos de adaptação do desenvolvedor, e não o inverso.
\end{description}
